\section{Introduction}

A reduced description can be excellent for passive dynamics yet still be insufficient for the interventions one intends to apply. Synergetics formalizes high-dimensional reduction through order parameters and slaving relations, while control and abstraction theory provide established criteria for deciding whether a reduced coordinate remains dynamically consistent under inputs.\cite{Haken1996,TabuadaPappas2005,GirardPappas2007} This article asks a deliberately narrower question: given a pre-declared classical synergetic macro/slaving map, is that particular map sufficient for a declared intervention family and response horizon? For the full retained-variable trajectory used here, the answer reduces to the standard controlled-projectability/closure condition; the condition itself is not claimed as new. The manuscript contribution is instead a restricted synergetics-centered diagnostic synthesis: connect the pre-declared map to that established criterion, test it in prospectively frozen minimal examples in two domains, retain successful countercontrols and negative results, and show in the physical example how bounded output-preserving hidden-state preparation can alter a later intervention response. The neural WEAK/FAIL/Gram/PARK evidence and the exact power-grid mean/COI countercontrol remain part of the main evidential chain rather than being hidden in supplementary material.

Synergetics formalizes the reduction of high-dimensional dynamics near organized regimes through order parameters and slaving relations: a small set of slow or unstable coordinates parametrizes degrees of freedom that relax more rapidly.\cite{Haken1996} Such reductions are naturally read as statements about passive dynamics, invariant or attracting manifolds, and the ability of a reduced description to reproduce selected trajectories. Control and abstraction theory ask a complementary question. If external inputs are admitted, when does a quotient or reduced coordinate evolve consistently under those inputs, and when do states identified by the reduction remain behaviorally indistinguishable? Controlled projectability, bisimulation, homomorphism, predictive-state representations, input-output computational mechanics, and causal abstraction all provide established languages for versions of that question.\cite{TabuadaPappas2005,Littman2001,Givan2003,RavindranBarto2003,BarnettCrutchfield2015,GirardPappas2007,Chalupka2015,Chalupka2016,Rubenstein2017,BeckersHalpern2019,LiKabaRavanbakhsh2025,XiaBareinboim2025,Geiger2025}

The central question is therefore deliberately restricted: \emph{given a pre-declared classical synergetic macro/slaving map, is that particular map sufficient for a declared intervention family and response horizon?} For the full retained-variable trajectory used here, this question reduces exactly to the established controlled-projectability/closure condition. The point is not to replace the abstraction literature. Rather, the condition supplies a precise diagnostic boundary for a synergetic reduction selected independently of the intervention-response test.

The contribution is a restricted synergetics-centered diagnostic synthesis. We connect a pre-declared slaving/order-parameter map to the established controlled-projectability criterion, test the resulting intervention-sufficiency diagnostic under prospectively frozen minimal examples in two domains, retain exact successful countercontrols, and show in the physical example how an output-preserving hidden-state preparation can alter the subsequent intervention response. We do not claim ownership of a generic state-equivalence criterion or a generic control method. The ingredients are individually close to or subsumed by established work; the contribution lies in their explicit synergetics-centered organization and in the frozen evidence chain.

Two features of that evidence chain are important. First, negative and limiting controls remain visible. In the neural branch, a learned two-dimensional response coordinate predicts held-out responses essentially exactly but is only WEAK relative to an equal-dimensional raw-parameter PCA baseline; a later nuisance-invariance gate is frozen as a specification-classification FAIL; and a symmetry-aware two-dimensional Gram representation is equally predictive and invariant. The learned-coordinate direction is therefore parked rather than promoted. In the power-grid branch, the representative-machine macro fails under a localized hidden-machine disturbance, but the arithmetic mean/center-of-inertia (COI) coordinate is exactly closed. The grid result is consequently coordinate-specific, not an argument against low-dimensional aggregation.

Second, the examples were prospectively specified before execution. Model classes, interventions, horizons, primary metrics, comparators, numerical methods, and decision thresholds were frozen before result inspection. Weak and failed outcomes were retained rather than repaired, and successful countercontrols were not omitted. This process discipline is part of the evidential interpretation, not a scientific claim by itself.

Section~II positions the manuscript against the closest parent literatures. Section~III states the diagnostic framework and the standard projectability boundary, including the frozen finite-horizon bridge. Sections~IV and V give the neural and power-grid witnesses with their mandatory negative and successful controls. Section~VI presents the output-preserving preparation benchmark in the same physical fibre. Section~VII synthesizes the evidence and limitations, and Sec.~VIII concludes with the restricted diagnostic question.

\venuefigure{figure_1_diagnostic_schematic}{Diagnostic schematic. Full-state fibre to a pre-declared synergetic macro/slaving map, frozen intervention family and retained response, and the standard controlled-projectability test, ending in either a sufficient macro or hidden-mode leakage. Conceptual only; no quantitative data.}{Diagnostic schematic; existing frozen SVG is referenced from the venue-neutral package.}

